\documentclass[ALCO,ThmDefs,Unicode,epreuves]{cedram}
\OneNumberAllTheorems

%\usepackage{amsmath,amsfonts,amsthm,amssymb}
%\usepackage{graphicx}
%\usepackage{fancyhdr}
%\usepackage{enumerate}
%\usepackage[justification=centering]{caption}
%\usepackage[labelformat=simple]{subcaption}

%----added by Nozaki---------
\usepackage{mathrsfs}
\usepackage{pdflscape}
\def\labelenumi {$(\theenumi)$}
\def\theenumii {\arabic{enumii}} 
\def\labelenumii {(\theenumii)} 


\providecommand{\bA}{\boldsymbol{A}}

\providecommand{\floor}[1]{\left\lfloor#1\right\rfloor}

\DeclareMathOperator{\ord}{ord}
\DeclareMathOperator{\tr}{tr}
\DeclareMathOperator{\GF}{GF}



%%%%%%%%%%%%%%%%%%%%%%%%%%%%%%%%%%%%%
\newcommand*{\mk}{\mkern -1mu}
\newcommand*{\Mk}{\mkern -2mu}
\newcommand*{\mK}{\mkern 1mu}
\newcommand*{\MK}{\mkern 2mu}

\hypersetup{urlcolor=purple, linkcolor=blue, citecolor=red}


\newcommand*{\romanenumi}{\renewcommand*{\theenumi}{\roman{enumi}}}
\newcommand*{\Romanenumi}{\renewcommand*{\theenumi}{\Roman{enumi}}}
\newcommand*{\alphenumi}{\renewcommand*{\theenumi}{\alph{enumi}}}
\newcommand*{\Alphenumi}{\renewcommand*{\theenumi}{\Alph{enumi}}}
%%%%%%%%%%%%%%%%%%%%%%%%%%%%%%%%%%



%%%%%%%%%%%%%%%%%%%%%%%%%%%%%%%%%%%%%%%%%%%
%%%%% Auteur

\author{\firstname{Sebastian} \middlename{M.} \lastname{Cioab\u{a}}} %%%%1

\address{University of Delaware\\ 
Department of Mathematical Sciences\\
Ewing Hall\\
Newark, DE 19716-2553, USA}

\email{cioaba@udel.edu}

\urladdr{sites.udel.edu/cioaba}

%\thanks{Sebastian M. Cioab\u{a} is supported by NSF grants DMS-1600768 and CIF-1815922.}

%%%%
\author{\firstname{Jack} \middlename{H.} \lastname{Koolen}} %%%%2

\address{School of Mathematical Sciences \\
University of Science and Technology of China\\
Wen-Tsun Wu Key Laboratory of the Chinese Academy of Sciences\\
 Hefei, Anhui, China}

\email{koolen@ustc.edu.cn}

\urladdr{http://staff.ustc.edu.cn/~koolen/}

%\thanks{Sebastian M. Cioab\u{a} is supported by NSF grants DMS-1600768 and CIF-1815922. Jack H. Koolen is partially supported by the National Natural Science Foundation of China (Nos.\ 11471009 and 11671376). He also acknowledges the financial support of the Chinese Academy of Sciences under its `100 talent' program.}

%%%%%
\author[\initial{H.} Nozaki]{\firstname{Hiroshi} \lastname{Nozaki}} %%%% 4

\address{Aichi University of Education\\
1 Hirosawa, Igaya-cho, Kariya\\
Aichi 448-8542, Japan}


\email{hnozaki@auecc.aichi-edu.ac.jp}

\urladdr{https://hnozaki.jimdo.com}

\thanks{Sebastian M. Cioab\u{a} is supported by NSF grants DMS-1600768 and CIF-1815922. Jack H. Koolen is partially supported by the National Natural Science Foundation of China (Nos.~11471009 and 11671376). He also acknowledges the financial support of the Chinese Academy of Sciences under its ``100 talent'' program. Hiroshi Nozaki is partially supported by JSPS Grants-in-Aid for Scientific Research Nos.~16K17569, 17K05155, and 18K03396.}

%%%%% Sujet

\keywords{second eigenvalue, bipartite regular graph, bipartite distance-regular graph, expander, 
linear programming bound. }

%\subjclass{10X99, 14A12, 11L05}
\subjclass{05B25, 05C35, 05C50, 05E30, 42C05, 94B65}

%%%%% Gestion
%
%\DOI{10.5802/alco.71}
%\datereceived{2018-05-02}
%\daterevised{2019-01-19}
%\datererevised{2019-03-04}
%\dateaccepted{2019-03-04}

%%%%% Titre et résumé

\title
[A spectral version of the Moore problem for bipartite regular graphs]
{A spectral version of the Moore problem for bipartite regular graphs}

\begin{abstract}
Let $b(k,\theta)$ be the maximum order of a connected bipartite $k$-regular graph whose second largest eigenvalue is at most $\theta$. In this paper, we obtain a general upper bound for $b(k,\theta)$ for any $0\leq \theta< 2\sqrt{k-1}$. Our bound gives the exact value of $b(k,\theta)$ whenever there exists a bipartite distance-regular graph of degree $k$, second largest eigenvalue $\theta$, diameter $d$ and girth $g$ such that $g\geq 2d-2$. For certain values of $d$, there are infinitely many such graphs of various valencies $k$. However, for $d=11$ or $d\geq 15$, we prove that there are no bipartite distance-regular graphs with $g\geq 2d-2$. 
\end{abstract}

\begin{document}

\maketitle

\section{Introduction}

Let $\Gamma=(V,E)$ be a connected $k$-regular simple graph with $n$ vertices. For $1\leq i\leq n$, let $\lambda_i(\Gamma)$ denote the $i$-th largest eigenvalue of the adjacency matrix of $\Gamma$. The eigenvalues have close relationships with other graph invariants. The smallest eigenvalue $\lambda_n(\Gamma)$ is related to the diameter, the chromatic number and the independence number (see~\cite[Chapter~4]{BH12} or~\cite{Ch89} for example). The second eigenvalue $\lambda_2(\Gamma)$ plays a fundamental role in the study of expanders~\cite{A86, AM85, BH12, HLW06}. Let $v(k,\theta)$ denote the maximum order of a connected $k$-regular graph $\Gamma$ with $\lambda_2(\Gamma) \leq \theta$. For $\theta < 2\sqrt{k-1}$, from work of Alon and Boppana, and Serre, we know that the value $v(k,\theta)$ is finite (see~\cite{CKNV16,MSS15}). In~\cite{CKNV16}, we obtained the following upper bound for $v(k,\theta)$. Let $\boldsymbol{T}(k,t,c)$ be the $t \times t$ tridiagonal matrix with lower diagonal $(1,1,\ldots,1,c)$, upper diagonal $(k,k-1,\ldots, k-1)$, and with constant row sum $k$. If $\theta$ is the second largest eigenvalue of $\boldsymbol{T}(k,t,c)$, then
\begin{equation}\label{vktau}
v(k,\theta) \leq 1+\sum_{i=0}^{t-3} k(k-1)^i+ \frac{k(k-1)^{t-2}}{c}. 
\end{equation}
Equality holds in~\eqref{vktau} if and only if there is a distance-regular graph of valency $k$ with second largest eigenvalue $\theta$, girth $g$ and diameter $d$ satisfying $g\geq 2d$. For $d>6$, there are no such graphs~\cite{DG81}. However, for smaller values of $d$, there are infinitely many values of $k$ and $\theta$ where the above inequality gives the exact value of $v(k,\theta)$. %Inequality~\eqref{vktau} is based on the linear programming bound for regular graphs~\cite{N} which will be improved in the case of bipartite graphs in Section~\ref{sec:lp}.

In this paper, we improve the above results from~\cite{CKNV16} for bipartite regular graphs. Let $b(k,\theta)$ denote the maximum order of a connected bipartite $k$-regular graph $\Gamma$ with $\lambda_2(\Gamma) \leq \theta$. Bipartite regular graphs $\Gamma$ with $\lambda_2(\Gamma) \leq \theta$ have been classified for $\theta=\sqrt{2}$~\cite{TY00}, $\theta=\sqrt{3}$~\cite{KS13}, and $\theta=2$~\cite{KS14}. We obtain a general upper bound for $b(k,\theta)$ for any $0\leq \theta< 2\sqrt{k-1}$. Our bound gives the exact value of $b(k,\theta)$ whenever there exists a bipartite distance-regular graph of degree $k$ with second largest eigenvalue $\theta$, diameter $d$ and girth $g$ such that $g\geq 2d-2$. For certain values of $d$, there are infinitely many such graphs of various valencies $k$. When $d=11$ or $d\geq 15$, we prove the non-existence of bipartite distance-regular graphs with $g\geq 2d-2$. Our results generalize previous work of H\o holdt and Justesen~\cite{HJ14} obtained in their study of graph codes and imply some results of Li and Sol\'{e}~\cite{CLS96} relating the second largest eigenvalue of a bipartite regular graph to its girth. The degree-diameter or Moore problem for graphs~\cite{MS} is about determining the largest graphs of given maximum degree and diameter. Given the connections between the diameter and the second largest eigenvalue of bipartite regular graphs (see~\cite{Ch89} for example), our Theorem~\ref{bound} can be interpreted as a spectral version of the Moore problem for bipartite regular graphs. 

In Section~\ref{sec:pre}, we describe some sequences of orthogonal polynomials and develop the preliminary results and notation that will be used in the paper. In Section~\ref{sec:lp}, we improve the linear programming bound from~\cite{N} for the class of bipartite regular graphs. In Section~\ref{sec:newbound}, we obtain the following upper bound for $b(k,\theta)$. Let $\boldsymbol{B}(k,t,c)$ be the $t \times t$ tridiagonal matrix with lower diagonal $(1,\ldots,1,c,k)$, upper diagonal $(k,k-1,\ldots, k-1,k-c)$, and constant row sum $k$. 
If $\theta$ is the second largest eigenvalue of $\boldsymbol{B}(k,t,c)$, then
\begin{equation}\label{bktau}
b(k,\theta) \leq 2\left(\sum_{i=0}^{t-4} (k-1)^{i}+ \frac{(k-1)^{t-3}}{c}+\frac{(k-1)^{t-2}}{c}\right). 
\end{equation}
We show that equality happens in~\eqref{bktau} when there is a bipartite distance-regular graph of degree $k$, second largest eigenvalue $\theta$ having $g\geq 2d-2$. Inequality~\eqref{bktau} generalizes some results of H\o holdt and Justesen~\cite{HJ14} (see Corollaries~\ref{coro:0} and~\ref{coro:1}), and of Li and Sol\'{e}~\cite{CLS96} (see Corollary~\ref{coro:Alon}). 
At the end of Section~\ref{sec:newbound}, we prove that the bound~\eqref{bktau} is better than~\eqref{vktau} for any $k$ and $\theta$. In Section~\ref{sec:bound_d}, we prove the non-existence of bipartite distance-regular graphs with $g\geq 2d-2$ for $d=11$ and $d\geq 15$. We conclude the paper with some remarks in Section~\ref{sec:conclusions}.



\section{Preliminaries} \label{sec:pre}
In this section, we describe some useful polynomials that will be used to prove our main result. For any integer $k\geq 2$, let $(F_i^{(k)})_{i\geq 0}$ be a sequence of orthogonal polynomials defined by the three-term recurrence relation: 
\begin{equation*}
F_0^{(k)}(x)=1, \qquad F_1^{(k)}(x)=x, \qquad F_2^{(k)}(x)=x^2 - k, 
\end{equation*}
and
\begin{equation} \label{eq:3-term}
F_i^{(k)}(x)=x F_{i-1}^{(k)}(x)- (k-1) F_{i-2}^{(k)}(x)
\end{equation}
for $i\geq 3$. 
The notation $F_i^{(k)}$ is abbreviated to $F_i$ for the rest of the paper. Let $q=\sqrt{k-1}$. The polynomials $(F_i)_{i\geq 0}$ form a sequence of orthogonal polynomials 
with respect to the positive weight
\[
w(x)=\frac{\sqrt{4q^2-x^2}}{k^2-x^2}
\]
on the interval $[-2q, 2q]$ 
(see~\cite[Section~4]{HOb}). 
The polynomials $F_i(q y)/ q^i$ in $y$ are called Geronimus
polynomials~\cite{G30,H15}. 
It follows from~\eqref{eq:3-term} that 
%\begin{align*}
%&F_0^{(k)}(x)=1,& &F_1^{(k)}(x)=x,& 
%&F_2^{(k)}(x)=x^2 - k,& \\
%&F_3^{(k)}(x)=x^3-(2k-1)x,& 
%&F_4^{(k)}(x)=x^4-(3k-2)x^2+k(k-1),&
%&%F_5^{(k)}(x)=x^5 -(4k-3)x^3+(3k-1)(k-1)x, 
%&
%\end{align*}
%and
\begin{equation} \label{eq:3-term2}
F_i(x)=(x^2-2k+2) F_{i-2}(x)- (k-1)^2 F_{i-4}(x)
\end{equation}
for $i\geq 5$. Note that for any $i\geq 0$, $F_{2i}(x)$ and $F_{2i+1}(x)$ are 
even and odd functions of $x$, respectively.

For $i\geq 0$, let $\mathscr{F}_{0,i}(x)=F_{2i}(\sqrt{x})$ and $\mathscr{F}_{1,i}(x)=F_{2i+1}(\sqrt{x})/\sqrt{x}$. 
It follows that $x^{\epsilon} \mathscr{F}_{\epsilon,i}(x^2)=F_{2i+\epsilon}(x)$ for $\epsilon\in \{0,1\}$. 
By~\eqref{eq:3-term2}, the polynomials $\mathscr{F}_{0,i}(x)$ and $\mathscr{F}_{1,i}(x)$ satisfy the following properties:
\begin{gather*}
\mathscr{F}_{0,0}(x)=1, \qquad 
\mathscr{F}_{0,1}(x)=x-k, \qquad
\mathscr{F}_{0,2}(x)=x^2-(3k-2)x+k(k-1), 
\\
\mathscr{F}_{1,0}(x)=1, \qquad 
\mathscr{F}_{1,1}(x)=x-(2k-1),
\end{gather*}
and 
\begin{equation} \label{eq:new_3-term}
\mathscr{F}_{\epsilon, i}(x)=(x-2k+2) \mathscr{F}_{\epsilon,i-1}(x)-(k-1)^2 \mathscr{F}_{\epsilon, i-2}(x)
\end{equation}
for any $i\geq 3$ if $\epsilon=0$, and $i\geq 2$ if $\epsilon=1$. 
Note that $k^{\epsilon} \mathscr{F}_{\epsilon,i}(k^2)=F_{2i+\epsilon}(k)=k(k-1)^{2i-1+\epsilon}=(k-1)^{2i-1+\epsilon}+(k-1)^{2i+\epsilon}$ for $2i+\epsilon \ne 0$. 
For $\epsilon\in \{0,1\}$, the polynomials $(\mathscr{F}_{\epsilon,i})_{i\geq 0}$ form a sequence of orthogonal polynomials with respect to the positive weight 
\[
w_\epsilon(x)=\frac{x^{\epsilon-1/2} \sqrt{4q^2-x}}{k^2-x}
\]
on the interval $[0, 4q^2]$. 

For $i\geq 0$, let $G_i(x)=\sum_{j=0}^{\floor{i/2}}F_{i-2j}(x)$. A simple calculation implies that
\begin{equation} \label{eq:G_i}
{G}_{i}(x)=
\frac{{F}_{i+2}(x)-(k-1)^2 {F}_{i}(x) }{x^2-k^2}
\end{equation}
for $i\geq 1$.
From Lemmas 3.3 and 3.5 in~\cite{CK07}, the polynomials $(G_i)_{i\geq 0}$ form a sequence of 
orthogonal polynomials with respect to 
the positive weight $(k^2-x^2)w(x)=\sqrt{4q^2-x^2}$ on the interval $[-2q,2q]$. From~\eqref{eq:3-term}, we deduce that
\[
G_i(x)=xG_{i-1}(x)-(k-1)G_{i-2}(x)
\]
for $i\geq 2$. 

Let $\mathscr{G}_{\epsilon, i}(x)$ denote the 
polynomial 
\begin{equation} \label{eq:def_scrG}
\mathscr{G}_{\epsilon, i}(x)=\sum_{j=0}^i
\mathscr{F}_{\epsilon, j}(x). 
\end{equation} 
It follows that 
$x^{\epsilon} \mathscr{G}_{\epsilon, j}(x^2)=G_{2j+\epsilon}(x)$. Using~\eqref{eq:G_i}, the polynomial $\mathscr{G}_{\epsilon, i}(x)$ can be expressed as
\begin{equation} \label{eq:pro_scrG}
\mathscr{G}_{\epsilon, i}(x)=
\frac{\mathscr{F}_{\epsilon, i+1}(x)-(k-1)^2\mathscr{F}_{\epsilon, i}(x) }{x-k^2}
\end{equation}
for any $i\geq 2$ if $\epsilon=0$, and
$i\geq 1$ if $\epsilon=1$. 
From Lemmas 3.3 and 3.5 in~\cite{CK07}, 
for $\epsilon\in \{0,1\}$, the polynomials $(\mathscr{G}_{\epsilon,i})_{i\geq 0}$ form a sequence of orthogonal polynomials with respect to 
the positive weight $(k^2-x)w_{\epsilon}(x)=x^{\epsilon-1/2}\sqrt{4q^2-x}$ on the interval $[0,4q^2]$. 

\begin{lemm} \label{lem:coe}
Let $p_l(i,j)$ be the coefficients in 
$x^{\epsilon} \mk \mathscr{F}_{\epsilon,i}(x)\mathscr{F}_{\epsilon,j}(x)\Mk =\Mk \sum_{l=0}^{i+j+\epsilon}\mk p_l(i,j)
\mathscr{F}_{0,l}(x)\mk $. Then
we have 
$p_0(i,j)=
k^{\epsilon}\mathscr{F}_{\epsilon,i}(k^2)\delta_{i,j}$, and 
$p_l(i,j)\geq 0$ for any $l,i,j$. Moreover 
$p_l(i,j)>0$ if and only if $|i-j|\leq l \leq i+j+\epsilon$. 
\end{lemm}
\begin{proof}
We have that
\[
F_{2i+\epsilon}(x)F_{2j+\epsilon}(x)\Mk =\Mk 
x^{2\epsilon}\mathscr{F}_{\epsilon,i}(x^2)\mathscr{F}_{\epsilon,j}(x^2)
\Mk =\Mk 
\sum_{l=0}^{i+j+\epsilon}p_l(i,j)
\mathscr{F}_{0,l}(x^2)
\Mk =\Mk 
\sum_{l=0}^{i+j+\epsilon}p_l(i,j)F_{2l}(x).
\]
By Theorem 3 in~\cite{N}, we obtain that 
$p_0(i,j)=F_{2i+\epsilon}(k) \delta_{i,j}=
k^{\epsilon}\mathscr{F}_{\epsilon,i}(k^2)\delta_{i,j}$, and 
$p_l(i,j)\geq 0$ for any $l,i,j$. Moreover 
$p_l(i,j)>0$ if and only if $|i-j|\leq l \leq i+j+\epsilon$. 
\end{proof}
Let $\Gamma$ be a connected regular bipartite graph. 
The adjacency matrix $\bA$ of $\Gamma$ can be expressed by 
\[
\boldsymbol{A}=
\begin{pmatrix}
\boldsymbol{O} & \boldsymbol{N}\\
\boldsymbol{N}^{\top} & \boldsymbol{O} 
\end{pmatrix}, 
\]
where $\boldsymbol{N}^{\top}$ is the transpose matrix of $\boldsymbol{N}$. 
The matrix $\boldsymbol{N}$ is called the \emph{biadjacency matrix} of $\Gamma$. 
It is not hard to see that 
\begin{equation} \label{eq:AB}
F_{2i}(\boldsymbol{A})= 
\begin{pmatrix}
\mathscr{F}_{0,i}(\boldsymbol{N}\boldsymbol{N}^{\top}) & \boldsymbol{O} \\
\boldsymbol{O} & \mathscr{F}_{0,i}(\boldsymbol{N}^{\top} \boldsymbol{N})
\end{pmatrix}. 
\end{equation}
%If $i$ is odd, then
%\[
%F_i(\boldsymbol{A})= 
%\begin{pmatrix}
%\boldsymbol{O} &\mathscr{F}_i(\boldsymbol{N}\boldsymbol{N}^{\top}) \boldsymbol{N} \\
% \mathscr{F}_i(\boldsymbol{N}^{\top} \boldsymbol{N})\boldsymbol{N}^{\top} & \boldsymbol{O} 
%\end{pmatrix}, 
%\]
%which is needed to write??
Since each entry of $F_{2i}(\boldsymbol{A})$ is non-negative~\cite{S66}, 
each entry of $\mathscr{F}_{0,i}(\boldsymbol{N}\boldsymbol{N}^{\top})$ is also non-negative. 
%From the positivity of $\mathscr{F}_{0,i}(\boldsymbol{N}\boldsymbol{N}^{\top})$, 
%we can obtain the linear programming bound for regular bipartite graphs. 


\section{Linear programming bound for bipartite regular graphs} 
\label{sec:lp}

In this section, we give a linear programming bound for bipartite regular graphs. For general regular graphs, a linear programming bound was obtained by Nozaki~\cite{N}. %As we will see in Table~\ref{tab:1} and Example~\ref{ex:1}, the {\em bipartite} linear programming bound below gives an improvement over the general linear programming bound from~\cite{N}. 
\begin{theo}
\label{thm:lp_bound}
Let $\Gamma$ be a connected bipartite $k$-regular graph with $v$ vertices. 
Let $\{\pm \tau_0, \ldots, \pm \tau_d \}$ be the set of distinct eigenvalues of $\Gamma$, 
where $\tau_0=k$. 
If there exists a polynomial $f(x)=\sum_{i= 0}^t f_i \mathscr{F}_{0,i}(x)$ such that 
$f(k^2)>0$, $f(\tau_i^2) \leq 0$ for each $i\in \{1,\ldots, d \}$, $f_0>0$, and $f_j \geq 0$ for each $j \in \{1,\ldots, t \}$, then 
\begin{equation} 
\label{eq:lp} 
v \leq \frac{2 f(k^2)}{f_0}. 
\end{equation}
Equality holds if and only if for each $i \in \{1,\ldots, d\}$, $f(\tau_i^2)=0$ and for each $j\in \{1,\dots,t\}$, $\tr (f_j\mathscr{F}_{0,j}(\boldsymbol{N}\boldsymbol{N}^{\top}))=0$, 
and $\tr (f_j\mathscr{F}_{0,j}(\boldsymbol{N}^{\top} \boldsymbol{N}))=0$, 
where $\boldsymbol{N}$ is the biadjacency matrix of $\Gamma$.
If equality holds and $f_j>0$ for each $j \in \{1,\ldots, t\}$, then the girth of $\Gamma$ is at least 
$2t+2$. 
\end{theo}
\begin{proof}
From the spectral decomposition $\boldsymbol{N}\boldsymbol{N}^{\top}= 
\sum_{i=0}^d \tau_i^2 \boldsymbol{E}_i$, we deduce that 
\begin{equation} \label{eq:pr1}
f(k^2) \boldsymbol{E}_0+ \sum_{i=1}^d f(\tau_i^2) \boldsymbol{E}_i
=f(\boldsymbol{N}\boldsymbol{N}^{\top})=\sum_{i= 0}^t f_i \mathscr{F}_{0,i}(\boldsymbol{N}\boldsymbol{N}^{\top})=
f_0 \boldsymbol{I} + \sum_{i= 1}^t f_i \mathscr{F}_{0,i}(\boldsymbol{N}\boldsymbol{N}^{\top}), 
\end{equation}
where $\boldsymbol{I}$ is the identity matrix, $\boldsymbol{E}_0=(2/v)\boldsymbol{J}$, and $\boldsymbol{J}$ is the all-ones matrix. 
Taking traces in both sides of~\eqref{eq:pr1}, we get that
\begin{multline*}
 f(k^2)=\tr ( f(k^2) \boldsymbol{E}_0 )
\geq \tr \left( f(k^2) \boldsymbol{E}_0+ \sum_{i=1}^d f(\tau_i^2) \boldsymbol{E}_i\right) \\
= 
\tr \left( f_0 \boldsymbol{I} + \sum_{j= 1}^t f_j\mathscr{F}_{0,j}(\boldsymbol{N}\boldsymbol{N}^{\top}) \right) 
\geq \tr ( f_0 \boldsymbol{I})
=\frac{v f_0}{2}. 
\end{multline*}
Therefore, $v\leq 2 f(k^2)/f_0$. 
By using $\mathscr{F}_{0,j}(\boldsymbol{N}^{\top} \boldsymbol{N})$, 
we can obtain the same bound as~\eqref{eq:lp}. 

If equality holds in~\eqref{eq:lp}, then for each $i \in \{1,\ldots, d\}$, $f(\tau_i^2)=0$ and for each $j\in \{1,\dots,t\}$, $\tr (f_j\mathscr{F}_{0,j}(\boldsymbol{N}\boldsymbol{N}^{\top}))=0$ and $\tr (f_j\mathscr{F}_{0,j}(\boldsymbol{N}^{\top} \boldsymbol{N}))=0$. 
 For the adjacency matrix $\boldsymbol{A}$, the $(u,v)$-entry of $F_{j}(\boldsymbol{A})$ is the number of non-backtracking walks of length $j$ from $u$ to $v$~\cite{S66}. 
Since~\eqref{eq:AB} and $f_j>0$ for each $j \in \{1,\ldots, t\}$, there is no non-backtracking walk of length $2j$ from $u$ to $v$ for each $j \in \{1,\ldots, t\}$. Since $\Gamma$ is bipartite, the girth of $\Gamma$ is at least $2t+2$. 
\end{proof}
%\begin{remark}
%We may normalize $f_0=1$ in Theorem~\ref{thm:lp_bound}. 
%\end{remark}
%\begin{remark} \label{rem:2}

%\end{remark}

\section{Upper bound for bipartite graphs with given second eigenvalue}
\label{sec:newbound}

In this section, we obtain an upper bound on $b(k,\theta)$ using the {\em bipartite} linear programming bound given by Theorem~\ref{thm:lp_bound}. Let $c>0$ be a real number and $t\geq 4$ be an integer. Let $\boldsymbol{B}(k,t,c)$ be the $t \times t$ tridiagonal 
matrix 
with lower diagonal $(1,\ldots,1,c,k)$, upper diagonal $(k,k-1,\ldots, k-1,k-c)$, and constant row sum $k$. 
Let
\[
\boldsymbol{B}(k,3,1)=\begin{pmatrix}
0&k&0 \\
1&0&k-1 \\
0&k& 0
\end{pmatrix}.
\]


\begin{theo} \label{bound}
If $\theta$ is the second largest eigenvalue of $\boldsymbol{B}(k,t,c)$, then
\[
b(k,\theta) \leq M(k,t,c)=2\left(\sum_{i=0}^{t-4} (k-1)^{i}+ \frac{(k-1)^{t-3}}{c}+\frac{(k-1)^{t-2}}{c}\right). 
\] 
Equality holds if and only if there exists a bipartite distance-regular graph whose quotient matrix with respect to the distance-partition from a vertex is $\boldsymbol{B}(k,t,c)$ for $1\leq c<k$ or $\boldsymbol{B}(k,t-1,1)$ for $c=k$. 
\end{theo} 

\begin{proof}
%Let $\boldsymbol{B}=boldsymbol{B}(k,t,c)$. 
We first calculate 
the characteristic polynomial of $\boldsymbol{B}(k,t,c)$. 
The polynomials $F_i$, $G_i$, $\mathscr{F}_i$, and $\mathscr{G}_i$ are defined in Section~\ref{sec:pre}. 
Note that $F_i(x)$ is the characteristic polynomial of the principal $i \times i$ matrix formed by the first 
$i$ rows and $i$ columns of $\boldsymbol{B}(k,t,1)$ for $t>i+1$. 
By this fact and equations~\eqref{eq:3-term2} and~\eqref{eq:G_i}, we can compute 
\begin{multline*}
|x \boldsymbol{I}-\boldsymbol{B}(k,t,c)|\\
\begin{aligned}
&={%\small 
\begin{vmatrix}
x & -k & & & & \\
-1 & x & -(k-1) & & & \\
 & \ddots & \ddots & \ddots & & \\
 & &-1 & x & -(k-1)& \\
 & & & -c & x & -(k-c) \\
 & & & & -k & x 
\end{vmatrix}}\\
&=k {%\small 
\begin{vmatrix}
x & -k & & & \\
-1 & x & -(k-1) & & \\
 & \ddots & \ddots & \ddots & \\
 & &-1 & x & 0 \\
 & & & -c & -(k-c) \\ 
\end{vmatrix}}+x{%\small 
\begin{vmatrix}
x & -k & & & \\
-1 & x & -(k-1) & & \\
 & \ddots & \ddots & \ddots & \\
 & &-1 & x & -(k-1) \\
 & & & -c & x 
\end{vmatrix}}\\
&=-k(k-c)F_{t-2}(x)+x(x F_{t-2}(x)-c(k-1)F_{t-3}(x))\\
&= c\big(F_{t-2}(x)-(k-1)^2 F_{t-4}(x)\big)+(x^2-k^2)F_{t-2}(x)\\
&= c\big((x^2-k^2)F_{t-4}(x)+F_{t-4}(x)-(k-1)^2 F_{t-6}(x)\big)+(x^2-k^2)F_{t-2}(x)\\
&= (x^2-k^2)\left(c\sum_{i=0}^{ \lfloor (t-4)/2 \rfloor}F_{t-4-2i}(x)+F_{t-2}(x)\right).
%\\
%&=(x^2-k^2)\big((c-1){G}_{t-4}(x)+{G}_{t-2}(x)\big)
\end{aligned}
\end{multline*}
Note that
\[
c\sum_{i=0}^{ \lfloor (t-4)/2 \rfloor}F_{t-4-2i}(x)+F_{t-2}(x)=(c-1){G}_{t-4}(x)+{G}_{t-2}(x). 
\]
Since the zeros of $\mathscr{G}_{\epsilon,s-1}$ and $\mathscr{G}_{\epsilon,s}$ interlace on $(0,4(k-1))$, each zero of $(c-1)\mathscr{G}_{\epsilon,s-1}+\mathscr{G}_{\epsilon,s}$ is simple and belongs to $(0,4(k-1))$ except for the smallest zero. 
For $c=k$ the smallest zero is equal to $0$ because 
$
(k-1)\mathscr{G}_{\epsilon,s-1}(0)+\mathscr{G}_{\epsilon,s}(0)=0
$ 
by~\eqref{eq:new_3-term} and~\eqref{eq:def_scrG}. 
For $c>k$, the smallest zero is negative. From $x^{\epsilon}(x^2-k^2)((c-1)\mathscr{G}_{\epsilon,s-1}(x^2)+\mathscr{G}_{\epsilon,s}(x^2))=(x^2-k^2)((c-1){G}_{2s-2+\epsilon}(x)+{G}_{2s+\epsilon}(x))$, 
each non-zero real eigenvalue of 
$\boldsymbol{B}(k,t,c)$ has multiplicity 1, 
and if $c>k$, then $\boldsymbol{B}(k,t,c)$ has 
imaginary eigenvalues. 

Let $\mathfrak{f}_1(x)$ be the polynomial 
\[
\mathfrak{f}_1(x)=\frac{((c-1){G}_{t-4}(x)+G_{t-2}(x))^2}{x^2-\theta^2}=\sum_{i=0}^{ t- 3} f_i \mathscr{F}_{0,i}(x^2).
\]
We show that $\mathfrak{f}_2(x)=\sum_{i=0}^{ t - 3} f_i \mathscr{F}_{0,i}(x)$ satisfies the condition of the linear programming bound from Theorem~\ref{thm:lp_bound} for bipartite graphs. 
Note that $\mathfrak{f}_2(k^2)=\mathfrak{f}_1(k)>0$, and $\mathfrak{f}_2(\lambda^2)=\mathfrak{f}_1(\lambda) \leq 0$ for each $\lambda \in [-\theta, \theta]$. 
It suffices to show that $f_i>0$ for each $i\in \{0,1,\ldots, t- 3\}$. 
 
The polynomial $\mathfrak{f}_1(x)$ can be expressed by 
\begin{align*}
\mathfrak{f}_1(x)&\Mk =\Mk 
\frac{(c-1){G}_{t-4}(x)+{G}_{t-2}(x)}{x^2-\theta^2}
\left(c\sum_{i=0}^{\lfloor t/2 \rfloor-2}{F}_{t-4-2i}(x)+{F}_{t-2}(x) \right)\\
&\Mk =\Mk x^{2\epsilon}
\frac{(c\Mk -\Mk 1)\mathscr{G}_{\epsilon,\lfloor t/2 \rfloor-2 }(x^2)\Mk +\Mk \mathscr{G}_{\epsilon,\lfloor t/2 \rfloor-1 }(x^2)}{x^2-\theta^2}
\Mk \left(\Mk c\sum_{i=0}^{\lfloor t/2 \rfloor-2}\Mk \mathscr{F}_{\epsilon,i}(x^2)\Mk +\Mk 
\mathscr{F}_{\epsilon,\lfloor t/2 \rfloor-1}(x^2)\Mk \right)\Mk , 
\end{align*}
where $\epsilon=0$ if $t$ is even, 
and $\epsilon=1$ if $t$ is odd. Thus, 
\[
\mathfrak{f}_2(x)=x^{\epsilon}
\frac{(c-1)\mathscr{G}_{\epsilon,\lfloor t/2 \rfloor-2 }(x)+\mathscr{G}_{\epsilon,\lfloor t/2 \rfloor-1 }(x)}{x-\theta^2}
\left(c\sum_{i=0}^{\lfloor t/2 \rfloor-2}\mathscr{F}_{\epsilon,i}(x)+
\mathscr{F}_{\epsilon,\lfloor t/2 \rfloor-1}(x)\right).
\]

By Proposition~3.2 in~\cite{CK07}, 
$g(x)=((c-1)\mathscr{G}_{\epsilon,\lfloor t/2 \rfloor-2}+\mathscr{G}_{\epsilon,\lfloor t/2 \rfloor-1})/(x-\theta^2)$ has positive coefficients 
in terms of $\mathscr{G}_{\epsilon,0},\mathscr{G}_{\epsilon,1},\ldots, \mathscr{G}_{\epsilon,\lfloor t/2 \rfloor-2}$. This implies that $g(x)$ has positive 
coefficients in terms of $\mathscr{F}_{\epsilon,0},\mathscr{F}_{\epsilon,1},\ldots, \mathscr{F}_{\epsilon,\lfloor t/2 \rfloor-2}$.
Therefore $f_i>0$ for each $i=\{0,1,\ldots, t-3 \}$ by 
Lemma~\ref{lem:coe}. 

The polynomial $g(x)$ can be expressed by $g(x)\Mk =\Mk \sum_{i=0}^{\lfloor t/2 \rfloor-2}g_i \mathscr{F}_{\epsilon,i}(x)$. 
By Lemma~\ref{lem:coe}, we have
\[
f_0=\sum_{i=0}^{\lfloor t/2 \rfloor-2} c k^{\epsilon} g_i \mathscr{F}_{\epsilon,i}(k^2)=c k^{\epsilon} g(k^2).
\]
By applying Theorem~\ref{thm:lp_bound} to the polynomial $\mathfrak{f}_2(x)$, we have
\begin{align*}
b(k,\theta) &\leq \frac{2\mathfrak{f}_2(k^2)}{f_0}=2k^\epsilon\left( \sum_{i=0}^{\lfloor t/2 \rfloor-2}\mathscr{F}_{\epsilon,i}(k^2)+\mathscr{F}_{\epsilon,\lfloor t/2 \rfloor-1}(k^2)/c \right) \\
&=2\left( \sum_{i=0}^{t-4} (k-1)^{i}+ \frac{(k-1)^{t-3}}{c}+
\frac{(k-1)^{t-2}}{c}
\right).
\end{align*}

By Theorem~\ref{thm:lp_bound}, the bipartite graph attaining the bound $M(k,t,c)$ has girth at least $2t-4$, and at most $t$ distinct eigenvalues. Since the diameter is at most $t-1$, the graph satisfies $g\geq 2d-2$, where $g$ is the girth and $d$ is the diameter. By $g\geq 2d-2$, the graph becomes a distance-regular graph~\cite[Theorem~4.4]{ADF16},~\cite{TY00}, and it must have the quotient matrix $\boldsymbol{B}(k,t,c)$ for $1\leq c < k$, or $\boldsymbol{B}(k,t-1,1)$ for $c=k$ (see Proposition~\ref{prop:c2} below). 
Conversely the distance-regular graph with the quotient matrix $\boldsymbol{B}(k,t,c)$ clearly attains the bound $M(k,t,c)$. 
\end{proof}
Note that $\Gamma$ is a distance-regular graph with the quotient matrix $\boldsymbol{B}(k,d+1,c)$ if and only if $\Gamma$ is a connected bipartite $k$-regular graph that has only $d+1$ distinct eigenvalues, and whose girth is at least $2d-2$. 
Table~\ref{tab:1} shows the known examples attaining the bound $M(k,d+1,c)$~\cite[Section~6.11]{BCN}. 
%Remark that the complete bipartite graph $K_{k,k}$ gives $b(k,0)=2k$. 

%%%%%%%%%%%%%%%%%%%%%%%%%%%%%%%%%%%%%%%%%%%%%%%%%%%%%%%%


\begin{exam} \label{ex:1}
Recall that $v(k, \theta)$ denotes the maximum order of a connected (not necessarily bipartite) $k$-regular graph whose second largest eigenvalue is at most $\theta$. We have $v(3,1)=10$, which is attained by the Petersen graph~\cite{CKNV16} and $b(3,1)=8$ from Table 1, which is attained by the bipartite incidence graph of the symmetric $(4,3,2)$-design. 
\end{exam}

The following is the bipartite version of Theorem 5 in~\cite{N}. 
\begin{coro}
Let $\Gamma$ be a bipartite distance-regular graph of order $n$ with quotient matrix $\boldsymbol{B}(k,t,c)$ with respect to the distance-partition from a vertex. Then $\lambda_2(\Gamma)\leq \lambda_2(\Gamma')$ for any bipartite $k$-regular graph $\Gamma'$ of order $n$. 
\end{coro}
\begin{proof}
Assume that there exists a graph $\Gamma'$ of order $n$ such that $\lambda_2(\Gamma')<\lambda_2(\Gamma)$. Then $\Gamma'$ also attains the bound from Theorem~\ref{bound}. This implies that $\Gamma'$ must have the eigenvalue $\lambda_2(\Gamma)$, which is a contradiction. 
\end{proof}

Let $\mu^{(j)}$ (\emph{resp}. $\lambda^{(j)}$) denote the largest zero of $F_j(x)$ (\emph{resp}. $G_j(x)$). 

\begin{prop}
For each $\theta\in [0,2\sqrt{k-1})$, there exist $t,c$ such that $\theta$ is the second largest eigenvalue of $\boldsymbol{B}(k,t,c)$.
\end{prop}
\begin{proof}
%We discuss a possible second eigenvalue $\lambda_2(t,c)$ of $\boldsymbol{B}(k,t,c)$. 

Note that $\lambda^{(j)}<\mu^{(j)}$ for $j\geq 1$ because 
$G_j(x)=\sum_{i=0}^{\lfloor j/2 \rfloor}F_{j-2i}(x) > 0$ for $x \geq \mu^{(j)}$.
The second eigenvalue $\lambda_2(t,c)$ of $\boldsymbol{B}(k,t,c)$ is equal to the largest zero of 
$(c-1)G_{t-4}(x)+G_{t-2}(x)$. 
 Since the zeros of $\mathscr{G}_{\epsilon,\floor{t/2}-2}$ and $\mathscr{G}_{\epsilon, \floor{t/2}-1}$ interlace, $\lambda_2(t,c)$ is a monotonically decreasing function in $c$. In particular, $\lim_{c\rightarrow \infty}\lambda_2(t,c)=\lambda^{(t-4)}$ with $t\geq 5$, $\lambda_2(t,1)=\lambda^{(t-2)}$, and 
$\lim_{c\rightarrow 0} \lambda_2(t,c)=\mu^{(t-2)}$. 
The largest zero $r^{(j)}$ of $G_j(x)+G_{j-1}(x)$ can be expressed by $r^{(j)}=2\sqrt{k-1} \cos \alpha$, where $\pi/(j+1) < \alpha < \pi/j$~\cite[Section~III.3]{BIb}. 
For $\lambda^{(j)}=2\sqrt{k-1} \cos \beta$, 
it follows from $r^{(j)}<\lambda^{(j)}$ that $\beta < \alpha < \pi/j$. This implies that
the possible value $\lambda_2(t,c)$ is between $\lim_{c\rightarrow k}\lambda_2(4,c)=0$ and $\lim_{t\rightarrow \infty} \lambda_2(t,1)=2\sqrt{k-1}$.
Therefore for each $\theta \in \mathopen[0 ,2\sqrt{k-1}\mathclose)$, there exist $t$, $c$ such that $\lambda$ is 
the second eigenvalue of $\boldsymbol{B}(k,t,c)$. 
\end{proof}


Note that for $\theta \in ( \lambda^{(t-2)}, \mu^{(t-2)}]$, $\theta$ is the second eigenvalue of both $\boldsymbol{B}(k,t,c_1)$ and $\boldsymbol{B}(k,t+2,c_2)$ for some $c_1$, $c_2$ with $0\leq c_1 < 1$, $c_2 > 0$. By the following proposition, we may assume $c\geq 1$ in Theorem~\ref{bound} to obtain better bounds. 

\begin{prop} \label{prop:c} 
Let $\theta\in \mathopen(\lambda^{(t-2)}, \mu^{(t-2)} \mathclose]$. 
Suppose $c_1$ and $c_2$ satisfy that $0\leq c_1 < 1$, $c_2 > 0$ and the second largest eigenvalues of $\boldsymbol{B}(k,t,c_1)$ and $\boldsymbol{B}(k,t+2,c_2)$ are $\theta$. 
Then we have $M(k,t,c_1)>M(k,t+2,c_2)$. 
\end{prop}
\begin{proof}
Since $(c_1-1)G_{t-4}(\theta)+G_{t-2}(\theta)=0$ holds, we have
 \[
c_1=-\frac{G_{t-2}(\theta)-G_{t-4}(\theta)}{G_{t-4}(\theta)}=-\frac{F_{t-2}(\theta)}{G_{t-4}(\theta)}. 
\]

\begin{landscape}
\begin{table}[htb]\centering
\vspace*{28mm}
\renewcommand{\arraystretch}{1.4}
%\renewcommand{\tabcolsep}{8pt}
\caption{Known bipartite graphs meeting the bound $M(k,d+1,c)$}\label{tab:1}
\begin{tabular}{|c|c|c|c|c|c|} 
\hline
$k$& $\theta$ & $b(k,\theta)$& $d$& $c$ &Name \\
\hline 
$2$&$2\cos(2\pi/n)$ & $n$ (even) & $n/2$& 1 & $n$-cycle $C_n$ \\
$k$&$0$ & $2k$& $2$ & $1$ & Complete bipartite graph $K_{k,k}$ \\
$k$& $\sqrt{k-\tau}$ & $2 (1+k(k-1)/\tau )$ & 3& $\tau$ & {%\small 
Symmetric $(v,k,\tau)$-design } \\ 
%$q+1$&$\sqrt{q}$ &$2(q^2+q+1) $&3 &1 & Incidence graph of $PG(2,q)$ &~\cite{S66,BCN}\\
{%\small 
$r^2-r+1$} &$r$ & 
{%\small 
$2(r^2+1)\times
%%%%%%%%%
(r^2-r+1)
%%%%%%%%%
$} &4 & 
{%\small 
$(r-1)^2$}
& {%\small 
$pg(r^2-r+1,r^2-r+1,(r-1)^2)$}\\
%& & {%\small 
%$(r^2-r+1)$} & & &\\
$q$&$\sqrt{q}$& $2q^2$ &4&$q-1$&$AG(2,q)$ minus a parallel class \\
$q+1$&$\sqrt{2q}$&$2\sum_{i=0}^3 q^i$ &4&1 & $GQ(q,q)$ \\
$q+1$&$\sqrt{3q}$ & $2\sum_{i=0}^5 q^i$ &6&1&$GH(q,q)$ \\
6 & 2& 162 & 4 & 2 & $pg(6,6,2)$ \\
\hline
\end{tabular}\\
\vrule height12pt width0pt $AG(2,q)$: affine plane, $GQ(q,q)$: generalized quadrangle, $GH(q,q)$: generalized hexagon, \\
$pg$: partial geometry, $q$: prime power, $r$: power of 2, \\
We use the bipartite incidence graph of an incidence structure. 
\end{table}
\end{landscape}

\noindent
Similarly $c_2\Mk =\Mk -F_{t}(\theta)/G_{t-2}(\theta)$ holds. 
By $\theta \Mk >\Mk  \lambda^{(t-2)}$, we have $F_{t-2}(\theta)\Mk =\Mk -c_1 G_{t-4}(\theta)<0$ and $F_{t}(\theta)=-c_2G_{t-2}(\theta)<0$. 
It therefore follows that 
\begin{align*}
M(k,t,c_1)-M(k,t+2,c_2) 
& =2k(k-1)^{t-2}\left(\frac{1}{c_1}-1-\frac{1}{c_2}(k-1)^2\right)\\
&= 2k(k-1)^{t-2}\left(-\frac{G_{t-4}(\theta)}{F_{t-2}(\theta)}-1+(k-1)^2\frac{G_{t-2}(\theta)}{F_{t}(\theta)}\right)\\
&= 2k(k-1)^{t-2}\left(-\frac{G_{t-2}(\theta)}{F_{t-2}(\theta)}+(k-1)^2\frac{G_{t-2}(\theta)}{F_{t}(\theta)}\right)\\
&=\frac{2k(k-1)^{t-2}G_{t-2}(\theta)}{F_{t-2}(\theta)F_{t}(\theta)}\big(-F_{t}(\theta)+(k-1)^2F_{t-2}(\theta) \big)\\
&=\frac{2k(k-1)^{t-2}(k^2-\theta^2)G_{t-2}(\theta)^2}{F_{t-2}(\theta)F_{t}(\theta)}>0. \qquad \qedhere
\end{align*}
\end{proof}
For $\theta \in ( \lambda^{(t-4)}, \lambda^{(t-3)}]$, $\theta$ is the second eigenvalue of both $\boldsymbol{B}(k,t,c_1)$ and $\boldsymbol{B}(k,t-1,c_2)$ for some $c_1$, $c_2$ with $c_1 \geq 1$, $c_2 \geq 1$. It follows that 
\begin{align*}
\frac{1}{c_1}&=-\frac{G_{t-4}(\theta)}{F_{t-2}(\theta)}= -\frac{G_{t-2}(\theta)}{F_{t-2}(\theta)}+1
=-\frac{\lambda G_{t-3}(\theta)-(k-1)G_{t-4}(\theta)}{F_{t-2}(\theta)}+1
\\
&=-\frac{\lambda G_{t-3}(\theta)}{F_{t-2}(\theta)}-\frac{k-1}{c_1}+1,
\end{align*}
and hence
\begin{equation} 
\label{eq:k/c}
\frac{k}{c_1}=-\frac{\lambda G_{t-3}(\theta)}{F_{t-2}(\theta)}+1
\end{equation}
for $t\geq 4$. Thus, 
if $\theta=\lambda^{(t-3)}$, then $c_1=k$. This implies that $k \leq c_1$. 
By the following proposition, we may assume $1 \leq c < k$ in Theorem~\ref{bound} to obtain better bounds. 

\begin{prop} \label{prop:c2}
Let $\theta\in \mathopen(\lambda^{(t-4)}, \lambda^{(t-3)} \mathclose]$. 
Suppose $c_1$ and $c_2$ satisfy that $k\leq c_1$, $c_2 \geq 1$ and the second largest eigenvalues of $\boldsymbol{B}(k,t,c_1)$ and $\boldsymbol{B}(k,t-1,c_2)$ are $\theta$. 
Then we have $M(k,t,c_1)\geq M(k,t-1,c_2)$. Moreover, equality holds if and only if 
$c_1=k$ and $c_2=1$. 
\end{prop}
\begin{proof}
From 
%$c_1=F_{t-2}/G_{t-4}$, 
$c_2=-F_{t-3}(\theta)/G_{t-5}(\theta)$ and $\eqref{eq:k/c}$, 
we have
\begin{align*}
\frac{M(k,t,c_1)- M(k,t-1,c_2)}{2(k-1)^{t-4}}&=
 1+\frac{k-1}{c_1}+\frac{(k-1)^2}{c_1}-\frac{1}{c_2}-\frac{k-1}{c_2}\\
&=1-\frac{1}{c_2}+(k-1)\left(
\frac{k}{c_1}-\frac{1}{c_2} \right)\\
&=1+\frac{G_{t-5}(\theta)}{F_{t-3}(\theta)}+(k-1)\left( 
-\frac{xG_{t-3}(\theta)}{F_{t-2}(\theta)}+1 +\frac{G_{t-5}(\theta)}{F_{t-3}(\theta)} \right)\\
&=\frac{G_{t-3}(\theta)}{F_{t-3}(\theta)}-(k-1)^2 \frac{F_{t-4}(\theta)G_{t-3}(\theta)}{F_{t-2}(\theta)F_{t-3}(\theta)}\\
&=\frac{G_{t-3}(\theta)}{F_{t-3}(\theta)F_{t-2}(\theta)}
(F_{t-2}(\theta)-(k-1)^2F_{t-4}(\theta))\\
&= (x^2-k^2)\frac{G_{t-3}(\theta)G_{t-4}(\theta)}{F_{t-3}(\theta)F_{t-2}(\theta)}\geq 0. 
\end{align*}
This implies the proposition. 
\end{proof}
The above results imply the following theorem. 
\begin{theo} \label{thm:lambda} 
Let $\lambda^{(j)}$ be the largest zero of $G_j(x)$ for $j\Mk \geq \Mk 1$. Then 
\[
\bigcup_{j= 1 }^\infty (\lambda^{(j)},\lambda^{(j+1)}]\Mk = (0,2 \sqrt{k-1}).
\]
 If $t\geq 4$ satisfies $\lambda^{(t-3)} < \theta \leq \lambda^{(t-2)}$, then 
\[
b(k,\theta) \leq M(k,t,c)=2\left(\sum_{i=0}^{t-4} (k-1)^{i}+ \frac{(k-1)^{t-3}}{c}+\frac{(k-1)^{t-2}}{c}\right), 
\] 
where $c=-F_{t-2}(\theta)/G_{t-4}(\theta)$. 
\end{theo}

The following results in~\cite{HJ12, HJ14} are obtained as corollaries of Theorem~\ref{thm:lambda}. 
\begin{coro}[{\cite{HJ12}}] \label{coro:0}
Let $\Gamma$ be a bipartite $n$-regular graph with $2m$ nodes. 
If 
$
\lambda_2(\Gamma) \leq \sqrt{n-1}
$,
then
\[
m \leq 1+\frac{n(n-1)}{n-\lambda_2^2(\Gamma)} 
\]
or, equivalently, 
\[
\lambda_2^2(\Gamma) \geq \sqrt{\frac{mn-n^2}{m-1}}. 
\]
\end{coro}	
\begin{proof}
This is immediate by Theorem~\ref{thm:lambda} for $t=4$. Indeed, 
since $\lambda^{(1)}$ is the largest zero of $G_1(x)=x$, we have
$\lambda^{(1)}=0$. 
Since $\lambda^{(2)}$ is the largest zero of $G_2(x)=x^2-(n-1)$, we have
$\lambda^{(2)}=\sqrt{n-1}$. 
Since $c=-F_{2}(\theta)/G_{0}(\theta)=n-\theta^2$, 
we have $M(n,4,c)/2=1+n(n-1)/(n-\theta^2)$. 
\end{proof}

\begin{coro}[{\cite[Theorem~4]{HJ14}}] \label{coro:1}
Let $\Gamma$ be a bipartite $n$-regular graph with $2m$ nodes. 
If 
$
\sqrt{n-1} \leq \lambda_2(\Gamma)\leq \sqrt{2(n-1)}
$, 
then
\[
m \leq n+\frac{n(n-1)}{2n-\lambda_2^2(\Gamma)-1}. 
\]
\end{coro}
\begin{proof}
This is immediate by Theorem~\ref{thm:lambda} for $t=5$. Indeed, 
since $\lambda^{(2)}$ is the largest zero of $G_2(x)=x^2-(n-1)$, we have
$\lambda^{(2)}=\sqrt{n-1}$.
Since $\lambda^{(3)}$ is the largest zero of $G_3(x)=x(x^2-2(n-1))$, we have
$\lambda^{(3)}=\sqrt{2(n-1)}$. 
Since 
$
c=-F_{3}(\theta)/G_{1}(\theta)=2n-\theta^2-1
$, 
we have 
\[
M(n,5,c)/2=n+n(n-1)/(2n-\theta^2-1).
\]
\end{proof}

For $0<\theta \leq \sqrt{k-1}$, the inequality $b(k,\theta) \leq 2(\theta^4+\theta^2+1)$ was obtained by Teranishi and Yasuno~\cite[Proposition~7.1]{TY00}. This bound is improved as follows. 
\begin{coro} \label{coro:2}
If $k^{1/4} <\theta \leq \sqrt{k-1}$, then
\[
b(k,\theta) \leq 2\left(1+\frac{k-1}{k-\theta^2} + \frac{(k-1)^2}{k-\theta^2}\right) < 2(\theta^4+\theta^2+1). 
\]
\end{coro}
\begin{proof}
Note that we have $\lambda^{(1)}=0$, $\lambda^{(2)}=\sqrt{k-1}$, and $c=-F_2(\theta)/G_0(\theta)=k-\theta^2$. By Theorem~\ref{thm:lambda}, for $0< \theta \leq \sqrt{k-1}$, we have
\[
b(k,\theta) \leq 2(1+(k-1)/(k-\theta^2) + (k-1)^2/(k-\theta^2)).
\] 
The inequality $1+\frac{k-1}{k-\theta^2} + \frac{(k-1)^2}{k-\theta^2} < \theta^4+\theta^2+1$ holds if and only if $k^{1/4}< \theta \leq \sqrt{k-1}$. The assertion therefore follows. 
\end{proof}

By Theorem~\ref{bound}, the following is immediate.
\begin{coro} \label{coro:Alon}
Let $\Gamma$ be a connected bipartite $k$-regular graph of order $v$. If $\theta$ is the second largest eigenvalue of $\boldsymbol{B}(k,t,c)$ and $v\geq M(k,t,c)$, then $\lambda_2(\Gamma) \geq \theta$ holds.
\end{coro}
Li and Sol\'{e}~\cite[Theorems~3 and~5]{CLS96} showed that if $\Gamma$
is of girth $g=2l$, then $\lambda_2(\Gamma) \geq 2\cos (\pi/l)$.
Corollary~\ref{coro:Alon} improves this result because we have $v \geq
M(k,l+1,1)$ when $g=2l$ and $\theta=2\cos (\pi/l)$ for $\boldsymbol{B}(k,l+1,1)$.


We prove that the bound~\eqref{bktau} is better than the bound $\eqref{vktau}$ for any $k$ and $\theta$. For $\eqref{vktau}$ we have a similar theorem to Theorem~\ref{thm:lambda}. For $j\geq 0$, denote $\mathcal{G}_j(x)=\sum_{i=0}^j F_i(x)$. 
\begin{theo}[\cite{CKNV16}] \label{thm:v2}
Let $r^{(j)}$ be the largest zero of $\mathcal{G}_j(x)$ for $j\geq 1$. Then $\bigcup_{j= 1 }^\infty (r^{(j)},r^{(j+1)}]= (-1,2 \sqrt{k-1})$. If $t\geq 3$ satisfies $r^{(t-2)} < \theta \leq r^{(t-1)}$, then 
\[
v(k,\theta) \leq N(k,t,c)=1+\sum_{i=0}^{t-3}k (k-1)^{i}+\frac{k(k-1)^{t-2}}{c}, 
\] 
where $c=-F_{t-1}(\theta)/\mathcal{G}_{t-2}(\theta)$. 
\end{theo} 

\begin{theo}
Let $k \geq 2$ be an integer. For $\theta \in (0,2 \sqrt{k-1})$, let $M(k,t_1,c_1)$ and $N(k,t_2,c_2)$ be defined as in Theorems~\ref{thm:lambda} and~\ref{thm:v2}, where $c_1=-F_{t_1-2}(\theta)/G_{t_1-4}(\theta)$, $c_2=-F_{t_2-1}(\theta)/\mathcal{G}_{t_2-2}(\theta)$, $\lambda^{(t_1-3)}<\theta \leq \lambda^{(t_1-2)}$, and $r^{(t_2-2)}<\theta \leq r^{(t_2-1)}$. Then 
\begin{equation*}
M(k,t_1,c_1) \leq N(k,t_2,c_2).
\end{equation*} 
Equality holds only if $t_1=t_2=t+1, \theta=\lambda^{(t-1)}, c_1=1$, and $c_2=k$. 
\end{theo}
\begin{proof}
Note that $\lambda^{(t-2)} < r^{(t-1)} < \lambda^{(t-1)}$ because $\mathcal{G}_{t-1}(x)=G_{t-1}(x)+G_{t-2}(x)$ for any $t\geq 3$.

Because $\theta \in (0,2\sqrt{k-1})=\cup_{j\geq 3}(\lambda^{(j-2)},\lambda^{(j-1)}]$, there is $t\geq 3$ such that $\theta\in (\lambda^{(t-2)},r^{(t-1)}]\cup (r^{(t-1)},\lambda^{(t-1)}]$. We consider each of the two possible cases $\lambda^{(t-2)}<\theta \leq r^{(t-1)}$ and $r^{(t-1)} < \theta \leq \lambda^{(t-1)} $ separately.

Suppose $\lambda^{(t-2)}<\theta \leq r^{(t-1)}$. Then $t_1=t+1$ and $t_2=t$. From Theorem~\ref{thm:v2}, a simple calculation yields that
\[
N(k,t_2,c_2)=2\sum_{i=0}^{t_2-3}(k-1)^i+(k-1)^{t_2-2}+\frac{k(k-1)^{t_2-2}}{c_2},
\]
and therefore, 
\begin{align*}
N(k,t,c_2)-M(k,t+1,c_1)&=\left(1+\frac{k}{c_2}-\frac{2k}{c_1} \right)(k-1)^{t-2}\\
&=\left(1-\frac{k\mathcal{G}_{t-2}(\theta)}{F_{t-1}(\theta)}+\frac{2kG_{t-3}(\theta)}{F_{t-1}(\theta)}\right)(k-1)^{t-2}\\
&=\left(1-\frac{kG_{t-2}(\theta)-kG_{t-3}(\theta)}{G_{t-1}(\theta)-G_{t-3}(\theta)}\right)(k-1)^{t-2}.
\end{align*}
Because the zeroes of $G_{t-2}$ and $G_{t-1}$ interlace, we get that $G_{t-1}(\theta)<0<G_{t-3}(\theta)$. Thus, $G_{t-1}(\theta)-G_{t-3}(\theta)<0$. If $G_{t-2}(\theta)-G_{t-3}(\theta)>0$, then it is clear that $N(k,t,c_2)>M(k,t+1,c_1)$. Otherwise, if 
$G_{t-2}(\theta)-G_{t-3}(\theta)<0$, then 
\begin{align*}
|kG_{t-2}(\theta)-\Mk kG_{t-3}(\theta)|\Mk -\Mk |G_{t-1}(\theta)\Mk -\Mk G_{t-3}(\theta)|
&\Mk =\Mk (k-1)G_{t-3}(\theta)\Mk +\Mk G_{t-1}(\theta)\Mk -\Mk kG_{t-2}(\theta)\\
&\Mk =\Mk (\theta - k)G_{t-2}(\theta)<0 
\end{align*}
which implies that $N(k,t_2,c_2)>M(k,t_1,c_1)$. 

Suppose $r^{(t-1)} <\theta \leq \lambda^{(t-1)}$. Then $t_1=t+1$ and $t_2=t+1$. Thus we have
\begin{multline*}
N(k,t+1,c_2)-M(k,t+1,c_1)=\left(k+1+\frac{k(k-1)}{c_2}-\frac{2k}{c_1}\right) (k-1)^{t-2}\\
\begin{aligned}
&=\left(k+1-\frac{k(k-1)\mathcal{G}_{t-1}(\theta)}{F_t(\theta)}+\frac{2kG_{t-3}(\theta)}{F_{t-1}(\theta)}\right) (k-1)^{t-2}\\
&=\left(k+1-\frac{k(k-1)(G_{t-1}(\theta)+G_{t-2}(\theta))}{F_t(\theta)}+\frac{2k(G_{t-1}(\theta)-F_{t-1}(\theta))}{F_{t-1}(\theta)}\right) (k-1)^{t-2}\\
&=\left(-\frac{k(k-1)G_{t-1}(\theta)}{F_t(\theta)}+\frac{2kG_{t-1}(\theta)}{F_{t-1}(\theta)}-(k-1)\left(1+k \frac{G_{t-2}(\theta)}{F_t(\theta)}\right) \right) (k-1)^{t-2}\\
&=\left(-\frac{k(k-1)G_{t-1}(\theta)}{F_t(\theta)}+\frac{2kG_{t-1}(\theta)}{F_{t-1}(\theta)}-\frac{(k-1)\theta G_{t-1}(\theta)}{F_t(\theta)} \right) (k-1)^{t-2}\\
&=\left( \frac{kF_t(\theta)-k(k-1)F_{t-1} (\theta)}{F_t(\theta) F_{t-1}(\theta)}+\frac{kF_t(\theta)-\theta(k-1)F_{t-1} (\theta)}{F_t(\theta) F_{t-1}(\theta)}\right) (k-1)^{t-2}G_{t-1}(\theta)\\
&=\left(\frac{k(\theta-k)\mathcal{G}_{t-1}(\theta)}{F_t(\theta) F_{t-1}(\theta)}+\frac{(\theta^2-k^2) G_{t-2}(\theta)}{F_t(\theta) F_{t-1}(\theta)} \right) (k-1)^{t-2}G_{t-1}(\theta)\\
&=\left(\frac{k\mathcal{G}_{t-1}(\theta)+(\theta +k) G_{t-2}(\theta)}{F_t(\theta) F_{t-1}(\theta)} \right) (k-1)^{t-2} (\theta-k) G_{t-1}(\theta)\geq 0. 
\end{aligned}
\end{multline*}
Equality holds only if $G_{t\Mk -\Mk 1}(\theta)\Mk = \Mk 0$ meaning that 
$\theta\Mk =\Mk \lambda^{(t-1)}$. Since $c_1\Mk =\Mk \frac{-F_{t-1}(\theta)}{G_{t-3}(\theta)}
\Mk =\Mk \frac{G_{t-3}(\theta)-G_{t-1}(\theta)}{G_{t-3}(\theta)}\Mk =1-\frac{G_{t-1}(\theta)}{G_{t-3}(\theta)}$ and $c_2=\frac{-F_{t}(\theta)}{\mathcal{G}_{t-1}(\theta)}=\frac{G_{t-2}(\theta)-G_{t}(\theta)}{G_{t-1}(\theta)+G_{t-2}(\theta)}=\frac{kG_{t-2}(\theta)}{G_{t-2}(\theta)}$, this means that
$c_1=1$, and $c_2=k$. 
\end{proof}

\section{Non-existence of certain distance-regular graphs}
\label{sec:bound_d}

In this section, we prove the non-existence of the graph that attains the bound in Theorem~\ref{bound} for $t>15$ and $t=12$.
Namely we prove the following. 
\begin{theo} \label{thm:non_existence}
Let $k$ and $c$ be two integers such that $k\geq 3$ and $1\leq c\leq k-1$. If $d=11$ or $d\geq 15$, there is no distance-regular graph $\Gamma$ with the quotient matrix $\boldsymbol{B} (k,d+1,c)$. 
\end{theo}
We prove Theorem~\ref{thm:non_existence} by the manner given by Fuglister~\cite{F87}.
Let $x=(t+1/t)\sqrt{k-1}$. The polynomial $G_i(x)$ can be expressed by 
\[
G_i(x)=\frac{\sqrt{(k-1)^i}}{t^i(t^2-1)}(t^{2i+2}-1). 
\]
The characteristic polynomial of $\boldsymbol{B}(k,d+1,c)$ is
$(x^2-k^2)\big( (c-1)G_{d-3}(x)+G_{d-1}(x) \big)$, and we have the expression 
\begin{align*}
\mathscr{S}_{d}(x)&=(c-1)G_{d-3}(x)+G_{d-1}(x)\\
&=\frac{\sqrt{(k-1)^{d-1}}}{t^{d-1}(t^2-1)}\left( 
t^{2d}+\frac{c-1}{k-1}t^{2d-2}- \frac{c-1}{k-1}t^2-1 \right).
\end{align*}
Let $\theta$ be an eigenvalue of $\Gamma$ that is not $\pm k$. Put $\theta=(\tau+1/\tau) \sqrt{k-1}$ for some complex number $\tau$. 
 Let $n$ be the order of $\Gamma$. The multiplicity $m_\theta$ is given by 
\[
m_\theta=\frac{nck(k-c)(k-1)^{d-2}}
{(k^2-\theta^2) \mathscr{S}_{d}'(\theta) f_{d-1}(\theta)},
\]
where $\mathscr{S}_{d}'(x)$ is the derivative with respect to $x$, $f_{d-1}=(x-1+c)G_{d-2}+(x-k+c)G_{d-3}-
(k-1)G_{d-4}$ and $G_{-1}=0$. 
From $\mathscr{S}_{d}(\theta)=0$, we can obtain 
\[
\tau^{2d-2}=\frac{(c-1) \tau^2+k-1}{(k-1)\tau^2+c-1}.
\]
Then we may calculate that
\[
f_{d-1}(\theta)= -\frac{c(k-c)\sqrt{(k-1)^{d-2}}}
{\tau^{d-2}((k-1)\tau^2+c-1)},
\]
and 
\[
\mathscr{S}_{d}'(\theta)\Mk =\Mk 
\frac{2 \sqrt{(k\Mk -\Mk 1)^{d\Mk -\Mk 4}}[(d\Mk -\Mk 1)(k\Mk -\Mk 1)(c\Mk -\Mk 1)(\tau^4\Mk +\Mk 1)\Mk + \Mk 
\big(d(k\Mk -\Mk 1)^2\Mk +\Mk (d\Mk -\Mk 2)(c\Mk -\Mk 1)^2\big)\tau^2]}
{\tau^{d-2}(\tau^2-1)^2\big((k-1)\tau^2+c-1\big)}. 
\]
 From these equations, the multiplicity $m_\theta$ can be expressed by 
\begin{align}
m_\theta&= \frac{nk(k-1)(\tau^2-1)^2 
\big( (c-1)\tau^2 + k-1 \big) \big((k-1)\tau^2+c-1 \big)}
{2(\theta^2-k^2) \tau^2 [(d-1)(k-1)(c-1)(\tau^4+1)+
 \big(d(k-1)^2+(d-2)(c-1)^2 \big) \tau^2
]} \nonumber 
\\
&=\frac{nk\big(\theta^2-4(k-1)\big)\big((c-1)\theta^2+(k-c)^2 \big)}
{2(\theta^2-k^2)[(d-1)(c-1)\theta^2+d(k-c)^2+2(c-1)(k-c)]} 
\nonumber 
\\
&=\frac{nk(k-1)\big(\phi-4\big)\big((c-1)(k-1)\phi+(k-c)^2 \big)}
{2\big((k-1)\phi-k^2\big)[(d-1)(c-1)(k-1)\phi+d(k-c)^2+2(c-1)(k-c)]}, \label{eq:quad} 
\end{align}
where $\theta^2=(k-1)\phi$. Unless $(k,c)=(2,1)$, expression~\eqref{eq:quad} gives a non-trivial rational quadratic polynomial in $\phi$. 

Set 
\[
H_d(x)=\frac{\mathscr{S}_d(x)}{x^\epsilon \sqrt{(k-1)^{d-3-\epsilon} }}, 
\]
where $\epsilon=1$ if $d$ is even, and $\epsilon=0$ if $d$ is odd. 
Let $z=x^2/(k-1)$. 
For $u=t^2$, we have $z=(t+1/t)^2=u+1/u+2$. We compute
\[
H_{2m+1-\epsilon}(z)=(c-1) P_{m-1,\epsilon}(z)+(k-1) P_{m,\epsilon}(z),
\]
where 
\[
P_{i,\epsilon}(z)=
\frac{ u^{2i+1-\epsilon}-1}
{u^{i-\epsilon}(u+1)^{\epsilon}(u-1)}. 
\] 
Note that $P_{i,\epsilon}(z)$ satisfy the recurrence relation
\[
P_{i,\epsilon}(z)=(z-2) P_{i-1,\epsilon}(z)-P_{i-2,\epsilon}(z)
\]
with the initial conditions $P_{0,\epsilon}(z)=1-\epsilon$, $P_{1,1}(z)=1$ and $P_{1,0}(z)=z-1$. This implies that 
$P_{i,\epsilon}(z)$ is a monic polynomial of degree $i$ 
with integer coefficients. 
Table~\ref{tb:prop_P} shows some useful identities 
involving polynomials $P_{i,\epsilon}(z)$. 
By~\eqref{eq:quad}, the polynomial $H_{d}(z)$ must split over the rationals into factors of degree at most $2$.

\begin{table}[htb] \centering
\caption{Identities involving $P_{i,\epsilon}(z)$} \label{tb:prop_P}
\renewcommand{\arraystretch}{1.3}
\begin{tabular}{|l|} 
\hline
$d=2m+1-\epsilon$ \\ 
\hline
$P_{m,\epsilon}(z)= (u^d-1)/u^{m-\epsilon} (u-1) (u+1)^\epsilon $\\
$P_{m-1,\epsilon}(z)= (u^{d-2}-1)/u^{m-1-\epsilon} (u-1) (u+1)^\epsilon$\\
$P_{m-1,\epsilon}(z)+P_{m,\epsilon}(z)=(u+1)^{1-\epsilon}(u^{d-1}-1)/u^{m-\epsilon} (u-1) $ \\
$-P_{m-1,\epsilon}(z)+P_{m,\epsilon}(z)= (u^{d-1}+1)/u^{m-\epsilon} (u+1)^\epsilon$ \\
\hline
\end{tabular}
\end{table}

\subsection{Case analysis modulo 2}

Let $c'=c-1$ and $k'=k-1$. If $c'$ and $k'$ have a factor in common, we may still factor out the content of $H_d(z)$. Call the resulting polynomial $\hat{H}_d(z)$. For $\hat{H}_d(z)$ modulo 2, there are three cases A--C, which are listed in Table~\ref{tb:mod2_1}. 
For natural numbers $n$ and $a$, let $\ord_n(a)$ be the non-negative integer $s$ such that $a=n^s b$ and 
$b$ is an integer that is not divisible by $n$. Suppose $\ord_n(0)= \infty$. 

\begin{table}[htb]\centering
 \caption{$\hat{H}_d(z)$ modulo 2, $w \in \{ 1,3,5\}$} \label{tb:mod2_1}
\renewcommand{\arraystretch}{1.3}
\begin{tabular}{|l|l|l|l|} \hline
Cases & Conditions& $\hat{H}_{d}(z) \pmod{2}$ & $d$ \\ \hline
A& $\ord_2(c')>\ord_2(k')$ & $P_{m,\epsilon}(z)$ &
$d=2^r w$ \\
B& $\ord_2(c')<\ord_2(k')$ & $P_{m-1,\epsilon}(z)$ &
$d-2=2^r w$ \\
C& $\ord_2(c')=\ord_2(k')$ & $P_{m-1,\epsilon}(z)+P_{m,\epsilon}(z)$ &
$d-1=2^r w$ \\
\hline
\end{tabular}
\end{table} 

Each root of $\hat{H}_d(z)$ is a root of one of the three irreducible polynomials of degree at most 2 over $\GF (2)$, which are listed in Table~\ref{tb:mod2_2}. There are also 
listed the results of the substitution $z=u+1/u+2$, as well as the multiplicative orders modulo 2 of the roots of the polynomials in $u$.

\begin{table}[htb]\centering
\caption{Irreducible polynomials over $\GF (2)$} \label{tb:mod2_2}
\renewcommand{\arraystretch}{1.3}
\begin{tabular}{|l|l|l|}
\hline 
$f(z)$ & $z=u+1/u+2$ & order of $u$ \\ \hline
$z$ & $(u+1)^2/u$ & 1\\
$z+1$ & $(u^2+u+1)/u$ & 3 \\
$z^2+z+1$ & $(u^4+u^3+u^2+u+1)/u^2$ & 5 \\ \hline
\end{tabular}
\end{table} 

If an expression $u^i-1$ occurs as a factor of $\hat{H}_d(z)$ modulo 2, then we must have $i=2^r w$ for $w\in \{1,3,5 \}$. 
From the identities in Table~\ref{tb:prop_P}, we can obtain 
the possible values for the diameter $d$ in Table~\ref{tb:mod2_1}. 

\subsection{Case analysis modulo 3}
 For $\hat{H}_d(z)$ modulo 3, there are three cases a--d, which are listed in Table~\ref{tb:mod3_1}. If $\ord_3(c')=\ord_3(k')=m$, then let $c''=c'/3^m$ and $k''=k'/3^m$. 


\begin{table}[htb]\centering
 \caption{$\hat{H}_d(z)$ modulo 3, $w \in \{1,2,4,5,8,10\}$} \label{tb:mod3_1}
\renewcommand{\arraystretch}{1.3}
\renewcommand{\tabcolsep}{3pt}
\begin{tabular}{|l|l|l|l|} \hline
Cases & Conditions& $\hat{H}_{d}(z) \pmod{2}$ & $d$ \\ \hline
a& $\ord_3(c')\Mk >\Mk \ord_3(k')$ & $\pm P_{m,\epsilon}(z)$ &
$d=3^r w$ \\
b& $\ord_3(c')\Mk <\Mk \ord_3(k')$ & $\pm P_{m-1,\epsilon}(z)$ &
$d-2=3^r w$ \\
c& $\ord_3(c')\Mk =\Mk \ord_3(k')$, $c''\Mk \equiv\Mk k'' \Mk \pmod{3} $ & $\pm (P_{m-1,\epsilon}(z)\Mk +\Mk P_{m,\epsilon}(z))$ &
$d-1=3^r w$ \\
d& $\ord_3(c')\Mk =\Mk \ord_3(k')$, $c''\Mk \equiv\Mk - k''\Mk \pmod{3} $ & $\pm (P_{m-1,\epsilon}(z)\Mk -\Mk P_{m,\epsilon}(z))$ &
$2d-2=3^r w$\\
\hline
\end{tabular}
\end{table} 


There are six irreducible polynomials of degree at most 2 over $\GF (3)$, which are listed in Table~\ref{tb:mod3_2}. 
We can obtain the possible values for the diameter $d$ in Table~\ref{tb:mod3_1} by a similar way to modulo 2. Here $w \in \{1,2,4,5,8,10 \}$.

\eject

\begin{table}[htb]
\centering
\renewcommand{\arraystretch}{1.3}
\caption{Irreducible polynomials over $\GF (3)$} \label{tb:mod3_2}
\begin{tabular}{|l|l|l|}
\hline 
$f(z)$ & $z=u+1/u+2$ & order of $u$ \\ \hline
$z-1$ & $(u-1)^2/u$ & 1 \\
$z$ & $(u+1)^2/u$ & 2\\
$z+1$ & $(u^2+1)/u$ & 4 \\
$z^2-z-1$ & $(u^2-u-1)(u^2+u+1)/u^2$ & 8 \\
$z^2+1$ & $(u^4+u^3+u^2+u+1)/u^2$ & 5 \\ 
$z^2+z-1$ & $(u^4-u^3+u^2-u+1)/u^2$ & 10 \\ 
\hline
\end{tabular}
\end{table} 

\subsection{A bound for the diameter}
Using a method similar to the one of Fuglister~\cite{F87}, we can 
obtain the possible values of $d$ in all cases A--C and a--d, which are listed in Table~\ref{tb:d_list}. 

\begin{table}[htb]\centering
\caption{Possible values of $d$} \label{tb:d_list}
\renewcommand{\arraystretch}{1.2}
\begin{tabular}{|l|l|l|}
\hline
\multicolumn{2}{|l|}{Case}  & \\
\cline{1-2}
(mod 2) & (mod 3) & Possible values of $d$\\
\hline
A & a & 3--6,8,10,12,24 \\
&b & 3--6,8,10,12,20,32 \\
&c,d & 3--6,10,16\\ 
B & a & 3--6,8,10,12,18,162 \\
&b & 3--8,10,12,14,26 \\
&c,d & 3--7,10,82\\
C & a& 3--6,9,81\\
&b& 3--7,11,17\\
&c&3--7,9,11,13,25\\
&d&3--7,13 \\
\hline
\end{tabular}
\end{table}

We eliminate several choices of $d$ from Table~\ref{tb:d_list} in this subsection. 
%\subsection{Elmination for small $d$}
\begin{prop}\label{prop:1}
There does not exist a 
distance-regular graph $\Gamma$ with the quotient matrix $\boldsymbol{B} (k,d+1,c)$ for $d=17,18,20,32,81,82,162$. 
\end{prop}
\begin{proof}
Using a computer, we can obtain the factorization of $c'P_{m-1,\epsilon}(z)+k'P_{m,\epsilon}(z)$ modulo $p$ into irreducible polynomials for given $d$, $k'$ and $c'$. 

For $d=18,81,82,162$, we can find an irreducible polynomial of degree at least $3$ as a factor of $c'P_{m-1,\epsilon}(z)+k'P_{m,\epsilon}(z)$ over $\GF (5)$ for each pair $(c',k') \in \GF (5)\times \GF (5)$. 

For $d=20,32$, we can find an irreducible polynomial of degree at least $3$ as a factor of $c'P_{m-1,\epsilon}(z)+k'P_{m,\epsilon}(z)$ over $\GF (7)$ for each pair $(c',k')\in \GF (7)\times \GF (7)$. 

For $d=17$, we can find an irreducible polynomial of degree at least $3$ as a factor of $c'P_{m-1,\epsilon}(z)+k'P_{m,\epsilon}(z)$ over $\GF (43)$ for each pair $(c',k')\in \GF (43)\times \GF (43)$.
\end{proof}

Bannai and Ito~\cite{BI79} proved the unimodal property of the multiplicities of the eigenvalues of 
Moore polygons. After this work, they also proved the rationality of the eigenvalues of Moore polygons~\cite{BI81}. 
The rationality of the eigenvalues is essential for the proof of the non-existence of Moore polygons~\cite{DG81}. 
In our case, the unimodal property of the multiplicities for the positive eigenvalues is easy. 
\begin{lemm} \label{lem:unimodal}
Let $\Gamma$ be a distance-regular graph with the quotient matrix $\boldsymbol{B} (k,d+1,c)$. Let $d'=\lfloor (d-1)/2 \rfloor$, which is the number of the positive non-trivial eigenvalues. 
Let $\theta_1, \ldots, \theta_{d'}$ be the positive non-trivial eigenvalues of $\Gamma$ with $\theta_1>\cdots > \theta_{d'}$. 
Let $m_{\theta_i}$ be the multiplicity of $\theta_i$. 
Then it follows that 
\[
m_{\theta_1} < m_{\theta_2} < \cdots < m_{\theta_{i-1}}< m_{\theta_i} \geq m_{\theta_{i+1}}>\cdots > m_{\theta_{d'}} \]
 for some $i \in \{1,\ldots, d'\}$. 
\end{lemm} 
\begin{proof}
The multiplicity $m_\theta$ of the eigenvalue $\theta$ can be expressed by equation~\eqref{eq:quad}. 
The function $m_\theta$ has no pole for $0< \phi < 4$ and $k\geq 3$. 
This implies the unimodal property of the multiplicities. 
\end{proof}

It is known that 
\[
G_i(x)= (k-1)^{i/2} U_i \left(\frac{x}{2 \sqrt{k-1}}\right), 
\]
where $U_i$ is the Chebyshev polynomial of degree $i$, which is defined by 
$U_i(\cos \theta)=\sin((i+1)\theta)/ \sin \theta$ (see~\cite{DSVb}). 
Thus the zeros of $G_i(x)$ are $2 \sqrt{k-1} \cos u_j^{(i)}$ for $j=1,\ldots, i$, where
$u_{j}^{(i)}=j \pi/(i+1)$. 
Since the expression 
\[
\mathscr{S}_{d}(x)=(c-1)G_{d-3}(x)+G_{d-1}(x)
=x G_{d-2}(x)-(k-c) G_{d-3}(x), 
\]
the positive zeros $\theta_{i}=2\sqrt{k-1} \cos \alpha_i$ of $\mathscr{S}_{d}(x)$ with $0<\alpha_1<\cdots < \alpha_{d'}< \pi /2$ 
satisfy 
$%(i-1) \pi /(d-2)
u^{(d-1)}_{i}<\alpha_i< 
%i \pi /(d-1)
u^{(d-2)}_{i}$ for each $i \in \{1,\ldots, d'\}$. 
%Let $\phi_i=\theta_i^2/(k-1)$. 
% i/(d)-(i-1)/(d-2)=id-2i-id+d/(d)(d-2)=d-2i/(d-1)(d-2)

Let $a(\phi)$ and $b(\phi)$ be the functions defined by
\begin{align*}
a(\phi)&= \frac{\phi -4}{(k-1) \phi - k^2}, \\
b(\phi)&= \frac{(c-1)(k-1)\phi+(k-c)^2}{(d-1)(c-1)(k-1)\phi+d(k-c)^2+2(c-1)(k-c)}\\
&=\frac{X_\phi}{(d-1)X_\phi+Y},
\end{align*} 
where $X_\phi=(c-1)(k-1)\phi+(k-c)^2$ and $Y=(k-c)(k+c-2)$. 
Note that $m_\theta=nk(k-1)a(\phi) b(\phi)/2$. 
Let $\theta_\alpha^2=(k-1) \alpha$ and 
$\theta_\beta^2=(k-1) \beta$. 
It follows that $m_{\theta_\beta}>m_{\theta_\alpha}$ if and only if 
\[
\frac{a(\beta)}{a(\alpha)}\cdot \frac{b(\beta)}{b(\alpha)}>1. 
\]
%%\frac{nk(k-1)\big(\phi-4\big)\big((c-1)(k-1)\phi+(k-c)^2 \big)}
%%{2\big((k-1)\phi-k^2\big)[(d-1)(c-1)(k-1)\phi+d(k-c)^2+2(c-1)(k-c)]}, \label{eq:quad} 
Let $\alpha=4 \cos^2 v$ and $\beta=4 \cos^2 w$. By direct calculation, 
\[
\frac{a(\beta)}{a(\alpha)}= 
1+ \frac{\cos 2v-\cos 2w}{1-\cos 2v} \cdot \frac{(k-2)^2}{(k-2)^2+2(k-1)(1-\cos 2w) } 
\]
and
\[
\frac{b(\beta)}{b(\alpha)}=1-\frac{Y(X_\alpha-X_\beta)} 
 {(d-1)X_\alpha X_\beta+Y X_\alpha}.
\]
%where $X_\alpha=(c-1)(k-1)\alpha+(k-c)^2$ and $X_\beta=(c-1)(k-1)\beta+(k-c)^2$. 

Let 
\[
A=\frac{\cos 2v-\cos 2w}{1-\cos 2v} \cdot \frac{(k-2)^2}{(k-2)^2+2(k-1)(1-\cos 2w) } 
\]
and 
\[
B=\frac{Y(X_\alpha-X_\beta)} 
 {(d-1)X_\alpha X_\beta+Y X_\alpha}.
\]
Note that 
\[
\frac{a(\beta)}{a(\alpha)}\cdot \frac{b(\beta)}{b(\alpha)}=(1+A)(1-B)=1+B\left( A\left(\frac{1}{B}-1\right)-1\right). 
\]
If $A(1/B-1)>1$ holds, then $a(\beta)b(\beta)/(a(\alpha)b(\alpha))>1$. 

\begin{lemm} \label{lem:L}
If $\pi /4 < v < w < \pi /2$ holds, then it follows that 
\[
A\left(\frac{1}{B}-1\right) >L(v,w)
\]
where $L(v,w):=\frac{3\sqrt{3} (d-1)}{4} \left(1-\frac{2}{k} \right)^2 
(1+\cos 2v) \sin^2 2 w$. 
\end{lemm}
\begin{proof}
We can calculate
\[
\frac{1}{B}-1=\frac{(d-1)X_\alpha X_\beta+Y X_\beta}{Y(X_\alpha-X_\beta)}. 
\]
Since it follows that 
\[X_\alpha%=(k-1)^2+(c-1)^2+2(c-1)(k-1) \cos 2v 
=\big((c-1) +(k-1) \cos 2v \big)^2+(k-1)^2(1-\cos^2 2v) \geq (k-1)^2\sin^2 2v,\] 
we have
\[
(d-1)X_\alpha X_\beta \geq (d-1)(k-1)^4\sin^2 2v\sin^2 2w
\]
and \[
YX_\beta \geq (k-c)(k+c-2)(k-1)^2\sin^2 2w.
\] 
It follows that 
\begin{align*}
Y(X_\alpha-X_\beta)&=2(k-c)(k+c-2)(c-1)(k-1)(\cos 2v-\cos 2w).
\end{align*}

For $1\leq c \leq k-1$, the function $f(c)=(k-c)(k+c-2)(c-1)$ is maximum at $c=1+(k-1)/\sqrt{3}$. It therefore follows that
\[
Y(X_\alpha-X_\beta)<\frac{4}{3\sqrt{3}}(k-1)^4(\cos 2v-\cos 2w). 
\]
Thus we obtain 
\begin{align*}
\frac{1}{B}-1&>
\frac{3\sqrt{3}\big( (d-1)(k-1)^2\sin^2 2v\sin^2 2w+(k-c)(k+c-2)\sin^2 2w\big)}{4(k-1)^2(\cos 2v-\cos 2w)}\\
&>\frac{3\sqrt{3} (d-1)\sin^2 2v\sin^2 2w}{4(\cos 2v-\cos 2w)},
%\\
%&=\frac{3\sqrt{3}\sin 2 w\big((d-1)(1-\cos^2 2v)+\frac{2k-3}{(k-1)^2}\big)}{8(w-v)}
\end{align*}
and hence
\begin{multline*}
A\left(\frac{1}{B}-1\right) \\
\begin{aligned}[b]
&> \frac{\cos 2v\Mk -\Mk \cos 2w}{1\Mk -\Mk \cos 2v} \cdot \frac{(k\Mk -\Mk 2)^2}{(k\Mk -\Mk 2)^2\Mk +\Mk 2(k\Mk -\Mk 1)(1 - \cos 2w) } \cdot
\frac{3\sqrt{3} (d - 1)\sin^2 2v\sin^2 2w}{4(\cos 2v\Mk -\Mk \cos 2w)}\\
%&>\frac{1}{ \tan v}\cdot \frac{(k-2)^2}{(k-2)^2+2(k-1)(1-\cos 2w) } \cdot
%\frac{3\sqrt{3} (d-1)(\sin^2 2v) \sin 2 w}{4}\\
&>\frac{3\sqrt{3} (d-1)}{4} \cdot \frac{(k-2)^2}{(k-2)^2+4(k-1) }\cdot 
(1+\cos 2v) \sin^2 2 w\\
&=\frac{3\sqrt{3} (d-1)}{4} \left(1-\frac{2}{k} \right)^2 
(1+\cos 2v) \sin^2 2 w .
\end{aligned}\qedhere\hspace*{-4mm}
\end{multline*}
\end{proof}
%\frac{\sqrt{3}}{4} \frac{(d-1)(k-1)^2+2k-3}{(k-1)^2(w-v) \sin 2v}. 

%If $d$ is even, then $d'=d/2-1$, 
%$(d'-2) \pi /(d-2)<\alpha_{d'-1}< (d'-1) \pi /(d-1)$,
%and
%$(i-1) \pi /(d-2)<\alpha_i< i \pi /(d-1)$

%We now obtain the following proposition. 
\begin{lemm} \label{lem:L>1}
Suppose $\alpha=4\cos^2 v=\theta_\alpha^2/(k-1)$, $\beta=4\cos^2 w=\theta_\beta^2/(k-1)$, and $\pi /4 < v < w < \pi /2$. 
If $L(v,w)\geq 1$, then $m_{\theta_\alpha}<m_{\theta_\beta}$. 
\end{lemm}
\begin{proof}
By Lemma~\ref{lem:L} and $B>0$, we have
\begin{align*}
\frac{m_{\theta_\beta}}{m_{\theta_\alpha}}&=\frac{a(\beta)}{a(\alpha)}\cdot \frac{b(\beta)}{b(\alpha)}=(1+A)(1-B)\\
&=1+B\left( A\left(\frac{1}{B}-1\right)-1\right)>1+B\big( L(v,w)-1\big)\geq 1. \qedhere
\end{align*}
\end{proof}

\begin{lemm}\label{lem:num_irr}
Let $\Gamma$ be a distance-regular graph with the quotient matrix $\boldsymbol{B} (k,d+1,c)$. 
Let $u_j^{(i)}=j \pi /(i+1)$ and $d'=\lfloor (d-1)/2 \rfloor$. 
If $L(u_{d'-j-1}^{(d-2)},u_{d'-j}^{(d-2)})\geq 1$ for some integer $j$ with $0 \leq j < d'-(d+3)/4$, 
then the number of positive eigenvalues $\theta$ of $\Gamma$ such that $\theta^2$ is irrational is less than or equal to $j$. 
\end{lemm}
\begin{proof}
The inequality $j < d'-(d+3)/4$ implies $u_{d'-j-1}^{(d-2)}>\pi/4$. 
If $\theta=2 \sqrt{k-1} \cos u$, we write $m_\theta=m(u)$. 
 %Let 
%$\theta_{d'-j-1}'=2\sqrt{k-1} \cos u_{d'-j-1}^{(d-2)}$
%and 
%$\theta_{d'-j}'=2\sqrt{k-1} \cos u_{d'-j}^{(d-2)}$. 
By Lemma~\ref{lem:L>1}, we have
$
m(u_{d'-j-1}^{(d-2)})<m(u_{d'-j}^{(d-2)})
$. 
Note that the eigenvalues $\theta_i=2 \sqrt{k-1} \cos \alpha_i$ satisfy 
 $u^{(d-1)}_{i}<\alpha_i< u^{(d-2)}_{i}$. 
By $
m(u_{d'-j-1}^{(d-2)})<m(u_{d'-j}^{(d-2)})
$ and the unimodal property of the function $m(u)$, we have $m(\alpha_{d'-j-1})<m(\alpha_{d'-j})$. 
Note that if $\theta^2$ is irrational, then $m_\theta=m_{\theta'}$ for some eigenvalue $\theta'$ with $\theta' \ne \theta$. 
The assertion therefore follows from $m(\alpha_{d'-j-1})<m(\alpha_{d'-j})$ and Lemma~\ref{lem:unimodal}. 
\end{proof}

\begin{prop} \label{prop:2}
There does not exist a 
distance-regular graph $\Gamma$ with the quotient matrix $\boldsymbol{B} (k,d+1,c)$ for $d=11, 16,24,25,26$. 
\end{prop}
\begin{proof}
By Lemma~\ref{lem:num_irr}, we can estimate the number of irrational square eigenvalues $(k-1)\phi=\theta^2$ by checking $L(u_{d'-j-1}^{(d-2)},u_{d'-j}^{(d-2)})>1$ with a computer. 

For $d=11$, the number of irrational $\phi$ is at most $1$ for $k\geq 5$. 
For $k=3,4$, we can find an irreducible polynomial of degree at least $3$ as a factor of $c'P_{m-1,\epsilon}(z)+k'P_{m,\epsilon}(z)$ over $\mathbb{Q}$ for each pair $(c',k')$ with $0\leq c'\leq k'-1$. 
We can find an irreducible polynomial of degree at least $3$ or two irreducible polynomials of degree $2$ as a factor of $c'P_{m-1,\epsilon}(z)+k'P_{m,\epsilon}(z)$ over $\GF (2)$ for each pair $(c',k') \in \GF (2)\times \GF (2)$. 

For $d=16$, the number of irrational $\phi$ is at most $2$ for $k\geq 3$. 
We can find an irreducible polynomial of degree at least $3$ or three irreducible polynomials of degree $2$ as a factor of $c'P_{m-1,\epsilon}(z)+k'P_{m,\epsilon}(z)$ over $\GF (3)$ for each pair $(c',k') \in \GF (3)\times \GF (3)$. 

For $d=25$, the number of irrational $\phi$ is at most $2$ for $k\geq 6$.
For $k=3,4,5$, we can find an irreducible polynomial of degree at least $3$ as a factor of $c'P_{m-1,\epsilon}(z)+k'P_{m,\epsilon}(z)$ over $\mathbb{Q}$ for each pair $(c',k')$ with $0\leq c'\leq k'-1$. 
We can find an irreducible polynomial of degree at least $3$ or three irreducible polynomials of degree $2$ as a factor of $c'P_{m-1,\epsilon}(z)+k'P_{m,\epsilon}(z)$ over $\GF (3)$ for each pair $(c',k') \in \GF (3)\times \GF (3)$. 

For $d=24, 26$, the number of irrational $\phi$ is at most $2$ for $k\geq 4$. 
For $k=3$, we can find an irreducible polynomial of degree at least $3$ as a factor of $c'P_{m-1,\epsilon}(z)+k'P_{m,\epsilon}(z)$ over $\mathbb{Q}$ for each $c'$ with $0\leq c'\leq k'-1$.
We can find an irreducible polynomial of degree at least $3$ or three irreducible polynomials of degree $2$ as a factor of $c'P_{m-1,\epsilon}(z)+k'P_{m,\epsilon}(z)$ over $\GF (3)$ for each pair $(c',k') \in \GF (3)\times \GF (3)$. 
\end{proof}

Theorem~\ref{thm:non_existence} follows
from Table~\ref{tb:d_list} and Propositions~\ref{prop:1},~\ref{prop:2}. 


\section{Conclusions}\label{sec:conclusions}

In this paper, we studied $b(k,\theta)$, the maximum number of vertices in bipartite regular graph of valency $k$ whose second largest eigenvalue is at most $\theta$. Our results extend previous work from~\cite{CKNV16, HJ12, HJ14, CLS96,TY00}. Our general bound for $b(k,\theta)$ is attained whenever there exists a bipartite distance-regular graph of valency $k$, second largest eigenvalue $\theta$, girth $g$ and diameter $d$ with $g\geq 2d-2$. For $d=3$ and $g \geq 4$ all the point-block incidence graphs of symmetric designs give equality in our bound so the situation is well-understood. For $d\geq 4$ we only have the Van Lint-Schrijver geometry besides the generalized polygons. We believe that for $d \geq 5$ the only examples must have $c=1$ and are generalized polygons.

 
\longthanks{The authors thank Tatsuro Ito for providing useful information relating to Section~\ref{sec:bound_d} and the two anonymous referees for their useful comments and suggestions.}

\nocite{*}
\bibliographystyle{amsplain-ac}
\bibliography{ALCO_Cioaba_162}

\end{document}


% First paragraph after a section is not indented. If there is text
% below the title before the first section, it should be unindented
% like this.
\noindent
The aim of this document is to be at the same time a manual and a
sample of a \verb|.tex| file prepared in the class
\verb|cedram-alco.cls| for submissions accepted for publication. The
class \verb|cedram-ALCO| is just a wrapper around \verb|amsart.cls|
(version 2) implementing the layout of the journal, therefore all
authors already at ease with the classes of the AMS should find it
easy to adapt their source to \verb|cedram-alco.cls|. As a
corollary, all packages whose name starts with \verb|ams| (like
\verb|amsmath| or \verb|amsfonts|) are already automatically loaded
and you don't need to call them explicitly in the
preamble\footnote{There is one exception: the package \texttt{amscd}
 is not loaded by default.}.

